\section{Proof of the lifted discretization theorem}
\label{sec:appendix-discretization}

We first record two elementary stability estimates for Bayesian
reweighting.

\begin{lemma}[Changing the base measure]
\label{lem:bayes-base-stability}
Let $\mu$ and $\nu$ be probability measures, let
$0\leq\ell\leq1$, and suppose
$\int\ell\,\mathrm d\mu\geq z$ and
$\int\ell\,\mathrm d\nu\geq z$ for some $z>0$.  Define
\[
  \mu^\ell=\frac{\ell\mu}{\int\ell\,\mathrm d\mu},
  \qquad
  \nu^\ell=\frac{\ell\nu}{\int\ell\,\mathrm d\nu}.
\]
There is a finite constant $C_z$, depending only on $z$, such that
\begin{equation}
  \dH(\mu^\ell,\nu^\ell)\leq C_z\dH(\mu,\nu).
  \label{eq:bayes-base-stability}
\end{equation}
\end{lemma}

\begin{proof}
Take a common dominating measure and denote the two densities by
$p,q$.  Cauchy--Schwarz gives
\[
  \int\abs{p-q}
  \leq
  \left(\int(\sqrt p-\sqrt q)^2\right)^{1/2}
  \left(\int(\sqrt p+\sqrt q)^2\right)^{1/2}
  \leq2\sqrt2\,\dH(\mu,\nu).
\]
Thus the two normalizing constants differ by at most the same
quantity.  Splitting the difference of the normalized square-root
densities into a density term and a normalizer term, and using that
$x\mapsto x^{-1/2}$ is Lipschitz on $[z,1]$, proves
\eqref{eq:bayes-base-stability}.
\end{proof}

\begin{lemma}[Changing the potential]
\label{lem:bayes-potential-stability}
Let $\nu$ be a probability measure and let $\Phi,\Psi\geq0$ satisfy
\[
  \int e^{-\Phi}\,\mathrm d\nu\geq z,\qquad
  \int e^{-\Psi}\,\mathrm d\nu\geq z
\]
for some $z>0$.  If $\Phi-\Psi\in L^2(\nu)$, then
\begin{equation}
  \dH\!\left(
    \frac{e^{-\Phi}\nu}{\int e^{-\Phi}\,\mathrm d\nu},
    \frac{e^{-\Psi}\nu}{\int e^{-\Psi}\,\mathrm d\nu}
  \right)
  \leq C_z\norm{\Phi-\Psi}_{L^2(\nu)} .
  \label{eq:bayes-potential-stability}
\end{equation}
\end{lemma}

\begin{proof}
Use
$\abs{e^{-x}-e^{-x'}}\leq\abs{x-x'}$ and
$\abs{e^{-x/2}-e^{-x'/2}}\leq\tfrac12\abs{x-x'}$ for
$x,x'\geq0$.  The normalizing constants are controlled in
$L^1(\nu)$, and the square-root likelihoods are controlled in
$L^2(\nu)$.  The same normalizer split as in
\Cref{lem:bayes-base-stability} gives the result.
\end{proof}

\begin{lemma}[Prior approximation]
\label{lem:prior-approximation}
Under \Cref{ass:drift-approximation},
\begin{equation}
  \dH(\muz,\mu_{0,N})\leq\frac12\sqrt{\varepsilon_N}.
  \label{eq:prior-approximation}
\end{equation}
\end{lemma}

\begin{proof}
The terminal-value map from path space to $\HH$ is measurable.  The
data-processing inequality and \eqref{eq:drift-energy} give
\[
  \KL(\muz\,\|\,\mu_{0,N})
  \leq
  \KL(\mathbb Q^\theta\,\|\,\mathbb Q_N^\theta)
  =\frac12\varepsilon_N.
\]
With the convention \eqref{eq:hellinger},
$2\dH^2(\mu,\nu)\leq\KL(\mu\,\|\,\nu)$.  This proves
\eqref{eq:prior-approximation}.
\end{proof}

\begin{lemma}[Likelihood projection error]
\label{lem:likelihood-projection}
For every $R<\infty$ there is a constant $C_R$, independent of $N$,
such that
\begin{equation}
  \sup_{\norm y\leq R}
  \norm{
    \Phim_{\sigobs}(\,\cdot\,;y)
    -\Phim_{\sigobs,N}(\,\cdot\,;y)
  }_{L^2(\mu_{0,N})}
  \leq C_R r_N.
  \label{eq:likelihood-projection-error}
\end{equation}
\end{lemma}

\begin{proof}
Write $U_T^N=L_N+H_N$, where
$L_N=P_NU_T^N$ and $H_N=(I-P_N)U_T^N$.  Since the drift in
\eqref{eq:lifted-prior-sde} depends only on $P_NU_t^N$ and takes
values in $\HN$, the high component is a stationary OU process
independent of the low component.  Hence
\[
  H_N\sim\mathcal N(0,(I-P_N)\mathcal C),\qquad
  \mathbb E\norm{H_N}^2=r_N^2,\qquad
  \mathbb E\norm{H_N}^4\leq3r_N^4.
\]
By \Cref{lem:sde-moments},
$\sup_N\mathbb E\norm{L_N}^4<\infty$.

For $u=l+h$, with $l=P_Nu$ and $h=(I-P_N)u$, expansion of the two
quadratic potentials gives, uniformly for $\norm y\leq R$,
\begin{equation}
  \abs{
    \Phim_{\sigobs}(u;y)-\Phim_{\sigobs}(l;y)
  }
  \leq C_R\norm h(1+\norm l+\norm h).
  \label{eq:potential-projection-pointwise}
\end{equation}
Square this inequality and use the independence of $L_N,H_N$, the
uniform fourth-moment bound, and the Gaussian estimates above.  Since
$r_N\leq(\operatorname{Tr}\mathcal C)^{1/2}$,
\[
  \mathbb E
  \left[
    \norm{H_N}^2(1+\norm{L_N}+\norm{H_N})^2
  \right]
  \leq C r_N^2.
\]
Taking square roots proves
\eqref{eq:likelihood-projection-error}.
\end{proof}

\begin{proof}[Proof of \Cref{thm:discrete}]
The moment bounds in \Cref{lem:sde-moments} are uniform in $N$.
Markov's inequality therefore gives a ball $B(0,r)$ that has mass at
least $1/2$ under every $\mu_{0,N}$ and positive mass under $\muz$.
On that ball, both $\Phim_{\sigobs}$ and
$\Phim_{\sigobs,N}$ are bounded uniformly for $\norm y\leq R$.
All normalizing constants used below consequently have a common
positive lower bound $z_R$.

Introduce the intermediate posterior
\[
  \frac{\mathrm d\bar\mu_N^y}{\mathrm d\mu_{0,N}}(u)
  =
  \frac{e^{-\Phim_{\sigobs}(u;y)}}
       {\int e^{-\Phim_{\sigobs}(v;y)}\,\mu_{0,N}(\mathrm dv)}.
\]
By the triangle inequality,
\[
  \dH(\muy,\muyN)
  \leq\dH(\muy,\bar\mu_N^y)
      +\dH(\bar\mu_N^y,\muyN).
\]
\Cref{lem:bayes-base-stability,lem:prior-approximation} give
\[
  \dH(\muy,\bar\mu_N^y)
  \leq C_R\sqrt{\varepsilon_N}.
\]
\Cref{lem:bayes-potential-stability,lem:likelihood-projection} give
\[
  \dH(\bar\mu_N^y,\muyN)\leq C_R r_N.
\]
Combining the estimates proves
\eqref{eq:discretization-bound}.
\end{proof}
